% TOP_PROBLEM: {"id": 2012, "title": "Erdős Problem #257", "category_id": 1, "category": {"id": 1, "name": "number_theory", "display_name": "Number Theory", "description": "Properties of integers, prime numbers, Diophantine equations.", "slug": "number-theory", "order_index": 1, "created_at": "2026-07-31T15:26:25.670Z"}, "status": "open"}
% FIRST_POSTED: null
% ENTRY_KIND: statement_only
% Imported from: unsolvedmath_additional_500.tex; UnsolvedMath ID 2012
% !TeX program = lualatex
% Compile twice: lualatex 2012.tex
% Self-contained: no external bibliography, images, or \input files.
% Numeric section labels and visible numbers are original UnsolvedMath IDs.
\documentclass[11pt,a4paper]{article}
\usepackage[margin=24mm,headheight=14pt]{geometry}
\usepackage{fontspec}
\setmainfont{Latin Modern Roman}
\setsansfont{Latin Modern Sans}
\setmonofont{Latin Modern Mono}
\usepackage{amsmath,amssymb,mathtools}
\usepackage{unicode-math}
\setmathfont{Latin Modern Math}
\usepackage{microtype}
\usepackage{enumitem,xurl,needspace}
\usepackage[dvipsnames]{xcolor}
\usepackage{hyperref}
\hypersetup{unicode=true,colorlinks=true,linkcolor=MidnightBlue,urlcolor=MidnightBlue,
 pdftitle={UnsolvedMath Problem 2012},pdfauthor={Source-annotated compilation},bookmarksnumbered=true}
\usepackage{bookmark,tocloft,titlesec,fancyhdr}
\setlength{\cftsecnumwidth}{4.2em}
\cftsetpnumwidth{2.5em}
\cftsetrmarg{4em}
\titleformat{\section}{\Large\bfseries\raggedright}{\thesection}{0.7em}{}
\pagestyle{fancy}\fancyhf{}
\fancyhead[L]{\small\sffamily Additional UnsolvedMath statements}
\fancyhead[R]{\small Original catalogue IDs}
\fancyfoot[C]{\thepage}
\setlength{\parindent}{0pt}
\setlength{\parskip}{5pt plus 1pt}
\setlength{\emergencystretch}{3em}
\setcounter{tocdepth}{1}\setcounter{secnumdepth}{1}
\setlist[itemize]{leftmargin=3em,itemsep=4pt,topsep=4pt,parsep=0pt}
\allowdisplaybreaks
\newcommand{\N}{\mathbb{N}}
\newcommand{\Z}{\mathbb{Z}}
\newcommand{\Q}{\mathbb{Q}}
\newcommand{\R}{\mathbb{R}}
\newcommand{\C}{\mathbb{C}}
\newcommand{\F}{\mathbb{F}}
\newcommand{\A}{\mathbb{A}}
\newcommand{\PP}{\mathbb{P}}
\newcommand{\E}{\mathbb{E}}
\newcommand{\eps}{\varepsilon}
\newcommand{\dd}{\,\mathrm d}
\newcommand{\sref}[2]{\hyperlink{source:#1:#2}{[S#2]}}
\newif\ifshowcatalogue
\showcataloguetrue
\newcommand{\cataloguescope}[1]{\ifshowcatalogue
 \paragraph{Statement recorded in the supplied source.}
 #1\fi}
\begin{document}
% UnsolvedMath ID 2012; source catalogue code EP-257

\setcounter{section}{2011}

\section{Irrationality of Arbitrary Mersenne Reciprocal Subseries}\label{2012}

{\small\sffamily UnsolvedMath ID 2012 · Number Theory}

\subsection{Definitions and mathematical statement}

\paragraph{Shared notation.}Unless a section says otherwise, $\N=\{1,2,\ldots\}$,
$\Z$ is the ring of integers, $\Q,\R,\C$ are the rational, real, and complex
fields, $[N]=\{1,\ldots,\lfloor N\rfloor\}$, and $\log$ is the natural
logarithm. For real functions of a parameter tending to infinity,
$f\ll g$ and $f=O(g)$ mean $|f|\leq Cg$ eventually for some fixed $C$;
$f\gg g$ reverses this comparison; $f=o(g)$ means $f/g\to0$;
$f\sim g$ means $f/g\to1$; and $f\asymp g$ means both order bounds.
A subscript on $O$ or $\ll$ specifies allowed constant dependence.
The expression $n^{o(1)}$ in an upper bound means growth smaller than
$n^\epsilon$ for each fixed $\epsilon>0$. Local definitions govern any
exception, finite-field convention, or use of the same symbol for another
object. An asymptotic claim is not automatically an effective algorithm.



The series is interpreted as the limit of its partial sums in increasing index order. Irrational means not in $\Q$; transcendental means not a root of any nonzero polynomial in $\Q[x]$. The latter is a stronger assertion. \sref{2012}{1} \sref{2012}{2}

\paragraph{Statement recorded in the supplied source.}
Let $A\subseteq \mathbb{N}$ be an infinite set. Is $ \sum_{n\in A}\frac{1}{2^n-1} $ irrational? \sref{2012}{1}

\subsection{Short English statement}

Must every infinite subseries of the reciprocals of numbers one below powers of two have an irrational sum?

\subsection{Sources}

{\small

\begin{itemize}

\item[\hypertarget{source:2012:1}{S1}] ULAM.AI, \emph{UnsolvedMath}, supplied \texttt{problems.json}, record 2012 (EP-257), statement and background. The embedded literature-review date is 2026-08-17; that is the archive’s review date, not a fresh certification. Dataset landing page: \url{https://huggingface.co/datasets/ulamai/UnsolvedMath}.

\item[\hypertarget{source:2012:2}{S2}] T. F. Bloom, \emph{Erdős Problems}, Problem 257, \url{https://www.erdosproblems.com/257}. Reference imported from the supplied record; the live problem page was not independently checked for this compilation.

\item[\hypertarget{source:2012:3}{S3}] V. Kovač and T. Tao, On several irrationality problems for Ahmes series, Acta Mathematica Hungarica 175 (2025), 572-608, \nolinkurl{doi:10.1007/s10474-025-01528-0.} \url{https://doi.org/10.1007/s10474-025-01528-0}. Bibliographic information imported from the supplied record.

\item[\hypertarget{source:2012:4}{S4}] T. Tao and J. Teräväinen, Quantitative correlations and some problems on prime factors of consecutive integers, arXiv:2512.01739 (2025). \url{https://arxiv.org/abs/2512.01739}. Bibliographic information imported from the supplied record.

\item[\hypertarget{source:2012:5}{S5}] Erd\H{o}s, P., On arithmetical properties of Lambert series. J. Indian Math. Soc. (N.S.) (1948), 63-66. Bibliographic information imported from the supplied record.

\item[\hypertarget{source:2012:6}{S6}] Erd\H{o}s, P., On the irrationality of certain series. Math. Student (1968), 222--226. Bibliographic information imported from the supplied record.

\end{itemize}

}

\label{end:2012}
\end{document}
